\documentclass[11pt,a4paper]{article}
\usepackage[UTF8,scheme=plain]{ctex}
\usepackage[margin=1.25cm,top=1.15cm,bottom=1.15cm]{geometry}
\usepackage{enumitem,booktabs,tabularx,array,xcolor}
\usepackage{titlesec}
\setlength{\parindent}{0pt}
\setlength{\parskip}{3pt}
\setlist{nosep,leftmargin=1.2em}
\titleformat{\section}{\large\bfseries\color{blue!55!black}}{}{0pt}{}
\titlespacing*{\section}{0pt}{5pt}{2pt}
\renewcommand{\arraystretch}{1.13}
\begin{document}
\begin{center}
{\LARGE\bfseries 在线代理模型研究：工作逻辑、结果与下一步}\\[-1pt]
{\small 刘润泽（Runze）\quad SDM374 Project-2\quad 2026年11月11日}
\end{center}
\vspace{-4pt}

\section{目标与实验逻辑}
目标是在逐轮获取 solver 真值的在线过程中，比较代理模型帮助寻找低 $R$ 权重 $W$ 的能力。五个模型从相同的 2048 条初始数据出发，在每轮共同的候选池上按预测均值最低选点；solver 返回真实 $R$ 后，只反馈给该模型自身并重训。这样保持候选选项一致、模型轨迹独立。在线寻优和预测误差分别评价，测试集题目与标签应在策略冻结前封存。

候选池经历三版：V1 每池 50 个全局 Dirichlet 点 + 50 个低值训练点附近扰动；V2 每池 60 个全局点 + 240 个围绕初始前 25 个低值点的多尺度扰动；V3 按润泽提出的离线漫游思路，围绕最佳初始权重 $p01969$，沿离线筛选方向及更广方向，以步长 0.01--0.05 生成每池 300 点。当前统一选点规则为最小预测均值；离线方向只是候选生成先验，真实优劣由 solver 验证。

\section{已完成实验和模型表现对比}
\footnotesize
V1 pilot 中 5 模型各运行 5 轮，仅得到 11 个不同选点，最佳值只从 115.765 微降至 115.759；100 点开发测试没有覆盖 $R\le115.68$。近优压力集 50 点中只有 7 点落入 $R\le120.765$，local-GP 在这 7 点的 MAE 最低（1.192），但标签已揭晓。

\vspace{1pt}
\begin{tabularx}{\linewidth}{@{}l>{\centering\arraybackslash}X>{\centering\arraybackslash}X>{\centering\arraybackslash}p{2.1cm}@{}}
\toprule
模型 & V2 最佳 $R$ 均值$\pm$SD（命中/3） & V3 最佳 $R$ 均值$\pm$SD（命中/3） & 拟合秒/轮 V2 / V3 \\
\midrule
local-GP & $113.820\pm0.534$（3/3） & $114.238\pm0.464$（3/3） & 56.0 / 48.0 \\
HistGB & $113.950\pm1.577$（2/3） & $113.915\pm0.090$（3/3） & 6.3 / 4.04 \\
XGBoost & $113.959\pm0.785$（3/3） & $114.290\pm0.827$（3/3） & 3.0 / 2.09 \\
boosted-consensus & $114.287\pm1.057$（3/3） & $114.434\pm0.511$（3/3） & 9.3 / 6.11 \\
ridge2 & $115.765\pm0.000$（0/3） & $114.163\pm0.350$（3/3） & 0.5 / 0.29 \\
\bottomrule
\end{tabularx}
\normalsize

\textbf{解读：}V2 多尺度池中 local-GP 均值最好、XGBoost 拟合快且表现接近；V3 漫游池中 HistGB 均值最低且最稳定，XGBoost 得到跨两版最低单点 $R=113.444$。ridge2 从 V2 的 0/3 命中变为 V3 的 3/3。每版内五模型共享相同候选池；V2 与 V3 改变了候选生成方式，跨版差异不能解释为纯模型因果差异。
\section{数据资产与验证边界}
去重 SQLite/CSV 库维护了 419 个不属于初始 2048 的唯一 solver 权重点（433 次观测，重复已折叠）：pilot 在线 11、pilot 开发测试 100、近优压力 50、V2 118、V3 140；与初始集重合为零。数据库和原始调用记录均按来源归档，后续可复用。Pilot 测试和压力测试都已揭晓标签，\textbf{不是盲测}；V4.1 sealed test 尚未生成，正式 20$\times$50 运行也尚未完成。V3 实际 solver 成本约 21.4 分钟；Local-GP 每轮拟合约 48 秒，树模型约 2--6 秒，ridge2 约 0.3 秒。

\section{结论与后续改进}
\begin{enumerate}
\item 继续验证离线方向和步长：增加 solver 覆盖，比较方向命中率、低 $R$ 分布，并测试多个低值锚点。
\item 固定候选池后再优化选点器：对比最小均值、LCB/EI 与探索配额；把代理模型质量和策略收益分开归因。动态移动候选中心另作端到端实验。
\item 将重复种子先扩至约 8 组，报告完整轨迹与均值/波动；冻结模型和规则后再制作独立 V4.1 测试集，一次性揭盲。
\item 控制成本：先小规模筛选方向/步长，只计算能影响模型选择的误差指标；暂不把 Spearman 作为主指标，单独记录 solver 与拟合耗时。
\end{enumerate}
\vfill
{\footnotesize 详细设计、逐轮日志、报告、候选池和原始调用记录见 \texttt{analysis/online\_model\_study\_v1/}；唯一数据表位于 \texttt{solver\_database/}。本报告数值来自已归档开发实验，结论限于所测试的候选生成方式与三个随机种子。}
\end{document}


